\documentclass[11pt,a4paper]{article}
\usepackage[T1]{fontenc}
\usepackage{lmodern,amsmath,amssymb,amsthm,mathtools,microtype}
\usepackage[margin=27mm]{geometry}
\usepackage[colorlinks=true,linkcolor=blue,urlcolor=blue]{hyperref}
\newtheorem{theorem}{Theorem}
\newtheorem{lemma}[theorem]{Lemma}
\newtheorem{corollary}[theorem]{Corollary}
\newcommand{\Q}{\mathbb Q}\newcommand{\C}{\mathbb C}
\DeclareMathOperator{\Res}{Res}\DeclareMathOperator{\Tr}{Tr}
\title{A divisorial obstruction on the actual multiplication surface}
\author{}\date{17 November 2026}
\begin{document}\maketitle

Let $K=\Q(\sqrt2)$ and let $E/K$ and $P\in E(K)$ be the fixed
elliptic curve and non-torsion source point. Fix an integer $n\ne0$.
On the abelian surface $A=\Res_{K/\Q}E$, write $H$ for the generic
point and impose the actual multiplication relations
\[
                       A_0=[n]P,\qquad U=[n]H.
\]
All formulas below are rational functions on the common source chart.
Put $f=1/x$, $z=-x/y$, and
\[
 r=\frac{z(U)}{z(H)},\qquad
 s=\frac{z(U-A_0)}{z(H-P)},
\]
\[
 (C_1,\ldots,C_6)=
 \bigl(f(H),f(U),\overline{f(H-P)},\overline{f(U-A_0)},
              r-\bar s,r+\bar s-2n\bigr).
\]
The two occurrences of conjugation are the actual two-slot operations
of the restriction of scalars; no trace is assigned independently.

For a rational function $t$ on $A$, define
\[
 F(t)=\frac{t(t+1)(2t+1)(3t+1)(4t+1)(5t+2)}{120},\quad
 w_0=\frac{F(t)}{4(t-1)^8},
\]
\[
 W=\prod_{i=1}^6\left(1+\frac{C_i^2}{F(t)}\right),
 \qquad V=\Tr(W^{-1}),\qquad
 \mathcal C=V^2+w(2-V).
\]
The physical weight $w$ may be either $w_0$ or
$4w_0/(1-w_0)^2$. Every original rational-function denominator is
retained; in the second case $1-w_0$ is retained as well.

\begin{theorem}\label{thm:obstruction}
Suppose $t\in\Q(A)$ is nonconstant. For any $c\in\Q(A)$, put
\[
 q=2+wc^2,\qquad
 \delta_c=(q^2-2q-1)^2+4q(q^2-1)V.
\]
None of the rational functions
\[
 \mathcal C,\qquad (1+24V)\mathcal C,\qquad
 (1+3(V-1)^2)\mathcal C,\qquad \delta_c\mathcal C
\]
is a square in $\C(A)$, on its nonempty guarded rational-function
domain. The same holds after multiplication by a nonzero constant,
or with a finite product of the displayed multiplier functions.
One may also include any rational multiplier in $(V,w)$ that is
regular and nonzero at $(0,0)$, such as $V+2$ or $1+w(2-V)$.
\end{theorem}

\begin{proof}
Base change to $\C$ and identify
\[
                          A_{\C}=E_1\times E_2.
\]
Write $H_1,H_2$ for the two generic elliptic points and $P_1,P_2$
for the two conjugate source points. Let $D$ be an arbitrary prime
divisor of this smooth projective surface, with valuation $v_D$.

If $v_D(t)<0$, then $v_D(F(t))=6v_D(t)$ is even. If
$v_D(F(t))$ is odd, it is therefore a positive integer $\nu$, and
the residue of $t$ is one of the six distinct roots of $F$.
None is one, so $v_D(t-1)=0$. Consequently
\[
                            v_D(w)=\nu>0
\]
for both allowed physical weights. In the second case $w_0$ has
positive valuation, so $1-w_0$ is a unit. This argument applies
to every prime divisor; regularity of $t$ on the whole surface
is neither assumed nor required.

Consider the first slot of $W$. Two of its tests are
\[
                          f_1(H_1),\qquad f_2(H_2-P_2).
\]
On an elliptic curve, the zero divisor of $1/x$ is $2[O]$.
Thus the positive-valuation divisors of these two tests are
supported on $\{O_1\}\times E_2$ and $E_1\times\{P_2\}$,
respectively. These divisors have no common irreducible component;
their intersection has codimension two. At least one of the two
tests therefore has nonpositive valuation at $D$.

For every test in this slot, including an identically zero test
with the convention $v_D(0)=+\infty$, parity gives
\[
 v_D\left(1+\frac{C_i^2}{F(t)}\right)
                     =\min\{0,2v_D(C_i)-\nu\}\le0.
\]
There is no cancellation, since $2v_D(C_i)-\nu$ is odd when
finite and cannot be zero. The test with nonpositive valuation
makes its factor valuation at most $-\nu$. Hence
$v_D(W_1^{-1})\ge\nu$.

For the second slot, use $f_2(H_2)$ and $f_1(H_1-P_1)$.
Their zero divisors likewise have no common component, so
$v_D(W_2^{-1})\ge\nu$. Summing gives $v_D(V)\ge\nu$,
and therefore
\[
                  v_D(\mathcal C)=\nu.
\]
Indeed, $V^2$ has valuation at least $2\nu$, while $w(2-V)$
has valuation exactly $\nu$.

The multipliers $1+24V$ and $1+3(V-1)^2$ have valuation zero,
with leading values one and four. Any rational function regular
and nonzero at $(V,w)=(0,0)$ is likewise a unit here.
For $\delta_c$, the case $c=0$
is the first of these. Otherwise $v_D(wc^2)$ is odd and nonzero.
If it is positive, $q$ has residue two and $\delta_c$ has residue
one. If it is negative, write $e=v_D(q)<0$ and $m=v_D(V)\ge\nu$.
The summands of $\delta_c$ have respective valuations
\[
                              4e,\qquad 3e+m.
\]
The first is strictly smaller, so $v_D(\delta_c)=4e$ is even.
Products of these multipliers also have even valuation. Thus any
square identity in the theorem forces $v_D(F(t))$ to be even
for every prime divisor $D$ on $A_{\C}$.

If $t$ depends nontrivially on $H_1$, restrict to the generic
fiber over $E_2$ and then extend the field $\C(E_2)$ to an algebraic
closure. The resulting nonconstant rational function on an elliptic
curve gives a finite morphism to $\mathbb P^1$ of some degree $m>0$.
A point of this generic fiber corresponds to a divisor of the surface
dominating $E_2$; separability in characteristic zero preserves its
valuation after the field extension. Hence every point over each of
the six distinct roots of $F$ has even ramification index. Each such
fiber has total multiplicity $m$ and at most $m/2$ distinct points.
Its contribution to the ramification divisor is therefore at least
$m/2$. The six fibers contribute at least $3m$.

The Riemann--Hurwitz formula~\cite{RH} gives total ramification $2m$ for a
degree-$m$ map from an elliptic curve to $\mathbb P^1$, a contradiction.
If $t$ is constant on the generic $E_1$ fiber, then
$t\in\C(E_2)$; its assumed nonconstancy gives the same contradiction
on $E_2$. This proves the theorem.
\end{proof}

\section*{Scope}

The theorem concerns a uniform rational-function construction on the
entire actual multiplication surface for one fixed nonzero integer
$n$. It allows the slope to depend arbitrarily and rationally on
$H$, with $U=[n]H$ and $A_0=[n]P$ imposed from the outset.
It also permits the adaptive source function $c$ to depend on that
slope and those actual points. In particular it covers a nonconstant
single-cover substitution $t=H_0^d+z(H)^2$ with constant $H_0$.

Constant $t$ is explicitly outside the theorem. The result does not
exclude individual rational witnesses, a point-dependent arithmetic
selection that is not a rational function, or witnesses on a proper
subvariety or a ramified cover of the multiplication surface. It
does not disprove either twenty-one-variable or nine-variable
realization. It excludes a proposed uniform formula even after the
actual multiplication relation has been imposed.

The accompanying checker
\nolinkurl{negative_results/computation/actual_multiplication_surface_v58/check.py}
verifies the exact contact polynomial, its six simple roots, the two
weight quotients at those roots, and the multiplier expansions.
The geometric nonexistence statement follows from the proof above;
it is not inferred from a finite search.

\begin{thebibliography}{9}
\bibitem{RH}The Stacks Project Authors, \emph{The Stacks Project},
Section 53.12, Riemann--Hurwitz, tag 0C1B,
\url{https://stacks.math.columbia.edu/tag/0C1B}.
\end{thebibliography}
\end{document}
